\documentclass[10pt,a4paper]{amsart}
\usepackage[margin=19mm]{geometry}
\usepackage{amsmath,amssymb,mathtools}
\usepackage[scr=rsfs,cal=euler]{mathalfa}
\usepackage{xfrac}
\usepackage[unicode,hidelinks,bookmarks=false]{hyperref}
\newcommand{\ie}{i.e.\ }
\newcommand{\cf}{cf.\ }
\newcommand{\resp}{respectively\ }
\newcommand{\Subsection}{\subsection}
\providecommand{\nobreakdash}{\nobreak\hbox{-}}
\setlength{\emergencystretch}{2em}
\multlinegap0pt
\usepackage{siunitx}
\usepackage{float}
\usepackage{xcolor}
\usepackage{mathtools}
\usepackage{amssymb}
\usepackage{bm}
\usepackage{tikz}
\newtheorem{assumption}{Assumption}
\newtheorem{proposition}{Proposition}
\title{\LARGE \bf Tilt as a Certified Resource: Preserving\\ Motor Wrench-Rate Authority on Articulated Multirotors}
\author{Giuseppe Silano, Martin Saska}
\date{Source: arXiv:2609.21580v1 (CC BY 4.0)}
\begin{document}
\maketitle

\noindent\emph{Source and licence.} \LARGE \bf Tilt as a Certified Resource: Preserving\\ Motor Wrench-Rate Authority on Articulated Multirotors. Version arXiv:2609.21580v1 (CC BY 4.0), distributed by arXiv under the Creative Commons Attribution 4.0 International licence, http://creativecommons.org/licenses/by/4.0/, as recorded in the arXiv licence metadata for this exact version and shown on the abstract page https://arxiv.org/abs/2609.21580v1. Full licence notice: \texttt{licenses/arxiv-2609.21580-CC-BY-4.0.txt}.\par\smallskip

\noindent\textit{Editorial scope.} Counted unit: the article's proposition prop:degeneracy (source0.tex lines 706--713). The article's complete original proof of it is delivered with it (source0.tex lines 715--726, one \texttt{proof} environment of 12 lines). No mathematics is written, repaired or re-derived: the statement and the proof are the article's own lines, in the article's own notation, and no retained line is substituted or re-flowed.  Retained uncounted as the source-local context the proof needs: lines 702--704.  Nothing was omitted from any retained span, so the package carries no omission record.  No retained line contains a citation, so the wrapper carries no bibliography and no bibliography entry.  No retained line contains a figure environment, an image-inclusion command or drawing code, so no image is delivered and the package declares no asset dependency.  Editorial formatting only, in addition: an AMS \texttt{amsart} preamble that restates the article's own declarations and macros used by the retained lines, namely the environments \texttt{assumption}, \texttt{proposition} on separate counters, together with the standard packages those lines need.  The retained environments print in this document as Assumption 1, Proposition 1 in order of appearance, whereas the article prints the same items with the numbering of its own full document.  The complete native source of the article as distributed in the arXiv e-print for this exact version is bundled unchanged as \texttt{sources/210974/source0.tex}.
\begin{assumption}\label{as:regular}
$v^\star$ is interior to $\{v: 0<v<v^{\mathrm{sat}}\}$, $J_v(v^\star)$ has full row rank $m$, and $v^\star$ is a regular local maximizer of $L(\cdot,\alpha)$ subject to $w(v,\alpha)=w^\star$, so a unique Lagrange multiplier $\lambda\in\mathbb{R}^m$ exists.
\end{assumption}

\begin{proposition}[Zero-sum authority degeneracy]\label{prop:degeneracy}
Under Assumption~\ref{as:regular}, $\nabla_v L(v^\star)=J_v(v^\star)^{\!\top}\lambda$, and the constant-wrench authority ascent available through rotor-speed redistribution,
\begin{align}
    \gamma \;=\; \max\big\{\nabla_v L^{\!\top}\xi \;:\; J_v\xi = 0,\ |\xi|\le\bar a\big\},
    \label{eq:gamma}
\end{align}
is zero. Here $\xi\in\mathbb{R}^n$ is a candidate spin-acceleration perturbation.
\end{proposition}

\begin{proof}
Interiority removes the box constraints; full row rank of $J_v(v^\star)$ is the linear independence constraint qualification under which $\lambda$ exists uniquely. The first-order necessary condition for the constrained maximizer is therefore $\nabla_v L(v^\star)=J_v(v^\star)^{\!\top}\lambda$. Any feasible $\xi$ in \eqref{eq:gamma} lies in $\ker J_v(v^\star)$, whence
%
\begin{align}
    \nabla_v L^{\!\top}\xi
    \;=\; \big(J_v^{\!\top}\lambda\big)^{\!\top}\xi
    \;=\; \lambda^{\!\top}\big(J_v\xi\big) \;=\; 0 .
    \label{eq:zerosum}
\end{align}
%
The linear objective therefore vanishes on the entire feasible set, which is nonempty since $\xi=0$ is admissible, so its maximum is zero.
\end{proof}

\end{document}
